\chapter*{ABSTRACT}
\thispagestyle{empty}

The project report titled "AirGlove" was developed by students Chicu Andrei, Cobzari Ion, Crudu Alexandra, Gurduza Mihai, and Moraru Patricia at the Technical University of Moldova. It is written in English and contains an introduction, 5 chapters (Problem Definition \& Analysis, Domain Analysis, Proposed Solution, Comparative Analysis, Target Audience \& Stakeholders), a list of figures, a list of tables, bibliography, and appendices.

This report examines the limitations of conventional mouse and touchpad interaction in scenarios where a stable surface is unavailable, uncomfortable, or impractical, and frames the core problem, goals, and constraints for a surface-independent pointing solution. Building on this, the domain analysis reviews inertial pointing and wearable interaction approaches from both research and industry, highlighting recurring issues such as drift, user-to-user variability, and the need for robust mapping strategies. The proposed solution then introduces the AirGlove prototype, which combines an ESP32 microcontroller with an MPU6050/MPU6500 IMU and functions as a Bluetooth HID device to improve portability and reduce reliance on USB dongles commonly used by commercial air mice. Implementation choices are motivated by the full mid-air control pipeline: sensor fusion, smoothing and filtering, dead zones, gain shaping, and clutching; emphasizing that stability and usability depend on end-to-end processing rather than raw orientation estimation alone. A comparative analysis (SWOT and structured comparison table) positions the prototype against commercial products and research prototypes, clarifying where it performs well (openness, customizability, wearable form factor) and where trade-offs remain, including calibration effort, precision expectations, power consumption, and drift. Finally, the target audience and stakeholders section connects these technical decisions to realistic usage contexts, especially presentation control and accessibility-oriented interaction, and outlines improvement directions aimed at increasing reliability, comfort, and robustness across different users and environments.

\textbf{Keywords: } wearable input device, air mouse, IMU, ESP32, Bluetooth HID, sensor fusion, accessibility, ergonomics